\IfFileExists{sourceserifpro.sty}{\usepackage{sourceserifpro}}{}
\AtBeginDocument{\frenchspacing}
\usepackage[flushmargin,bottom]{footmisc}
\usepackage{xcolor}
\usepackage{pagecolor}
\definecolor{wfmTeal}{HTML}{56949F}
\definecolor{wfmSalmon}{HTML}{D7827E}
\definecolor{wfmYellow}{HTML}{EA9D34}
\definecolor{wfmPink}{HTML}{B4637A}
\definecolor{wfmLime}{HTML}{8B9A5A}
\definecolor{wfmBlue}{HTML}{286983}
\definecolor{wfmGreen}{HTML}{6F8F72}
\definecolor{wfmPurple}{HTML}{907AA9}
\definecolor{wfmText}{HTML}{575279}
\definecolor{wfmPageBg}{HTML}{FEFCFA}
\AtBeginDocument{\pagecolor{wfmPageBg}\color{black}}
\usepackage{titling}
\setlength{\droptitle}{0cm}
\usepackage{titlesec}
\titleformat*{\section}{\centering\large\bfseries\color{black}}
% Second-level headings sit one point above the 11pt body text.
\titleformat*{\subsection}{\fontsize{12pt}{14pt}\selectfont\bfseries\color{black}}
% Fourth-level headings: never numbered, bold, on their own line, with
% almost no space before the paragraph that follows.
\titleformat{\paragraph}[block]{\normalfont\normalsize\bfseries\color{black}}{}{0pt}{}
\titlespacing*{\paragraph}{0pt}{3ex plus 1ex minus .2ex}{0.2ex}
% Fifth-level headings: the same, in italic rather than bold.
\titleformat{\subparagraph}[block]{\normalfont\normalsize\itshape\color{black}}{}{0pt}{}
\titlespacing*{\subparagraph}{0pt}{3ex plus 1ex minus .2ex}{0.2ex}
% Sixth-level headings are run in by heading-six.lua as "Heading.---Text".
\usepackage{float}
\usepackage{fancyhdr}
\usepackage{etoolbox}
% Keep wrapped footnote text aligned, with a gap after each hanging marker.
\makeatletter
\patchcmd{\@makefntext}
  {\llap{\@makefnmark}}
  {\llap{\@makefnmark\hspace{0.5em}}}
  {}{\PackageError{wfmR}{Unable to set footnote marker spacing}{}}
\makeatother
\usepackage{flafter}
\IfFileExists{placeins.sty}{\usepackage{placeins}}{\newcommand{\FloatBarrier}{}}
\usepackage{caption}
\captionsetup[figure]{font=footnotesize, labelfont=bf}
\renewcommand{\figurename}{Chart}
\makeatletter
\def\fps@figure{tb}
\def\fps@table{tb}
\makeatother
\setcounter{topnumber}{3}
\setcounter{bottomnumber}{2}
\setcounter{totalnumber}{4}
\renewcommand{\topfraction}{0.85}
\renewcommand{\bottomfraction}{0.75}
\renewcommand{\textfraction}{0.12}
\renewcommand{\floatpagefraction}{0.95}
\newcommand{\wfmFooterMark}{\hspace{0.45em}\raisebox{0.05em}{\textcolor{wfmSalmon}{\rule{0.55em}{0.55em}}}}
\pagestyle{fancy}
\fancyhf{}
\fancyfoot[R]{\small \thepage\wfmFooterMark}
\renewcommand{\headrulewidth}{0pt}
\renewcommand{\footrulewidth}{0pt}
\fancypagestyle{plain}{%
  \fancyhf{}%
  \fancyfoot[R]{\small \thepage\wfmFooterMark}%
  \renewcommand{\headrulewidth}{0pt}%
  \renewcommand{\footrulewidth}{0pt}%
}
\AtBeginEnvironment{CSLReferences}{\small}
% Tables set a step smaller than body text: a table is scanned, not
% read line by line, so it takes tighter setting than prose, and the
% smaller size lets wide tables fit the measure without scaling.
% longtable covers Quarto's default kable output; tabular and
% tabularx cover the rest.
\AtBeginEnvironment{longtable}{\footnotesize}
\AtBeginEnvironment{tabular}{\footnotesize}
\AtBeginEnvironment{tabularx}{\footnotesize}
\captionsetup[table]{font=small}
